\SourceLeaf{35}
These investigations have achieved the aim of this section. We add a few more propositions based on similar principles, which we will often need later.

\Needspace{6\baselineskip}
\begin{center}\emph{Various theorems.}\end{center}

\Article{38}
\textsc{Problem.} \emph{Find how many positive numbers are smaller than a given positive number $A$ and coprime to it.}

For brevity, let the symbol $\phi$, placed before a given number, denote the number of positive numbers coprime to it and smaller than it. We therefore seek $\phi A$.

I. When $A$ is prime, clearly all numbers from $1$ to $A-1$ are coprime to $A$. So in this case
\[\phi A=A-1\]

II. When $A$ is a power of a prime, say $=p^m$, none of the numbers divisible by $p$ are coprime to $A$, whereas all the others are. So, from the $p^m-1$ numbers, we must discard $p,2p,3p....(p^{m-1}-1)p$. This leaves $p^m-1-(p^{m-1}-1)$, or $p^{m-1}(p-1)$. Thus
\[\phi p^m=p^{m-1}(p-1)\]

III. The remaining cases are easily reduced to these by the following proposition: \emph{If $A$ is resolved into pairwise coprime factors $M,N,P$ etc., then}
\[\phi A=\phi M.\,\phi N.\,\phi P\ \textit{etc.}\]
The proof is as follows. Let the numbers coprime to $M$ and smaller than $M$ be $m,m'.m''$ etc., so their number is $=\phi M$. Similarly, let the numbers coprime to $N,P$ etc., respectively, and smaller than them be $n,n',n''$ etc., $p,p',p''$ etc., whose numbers are $\phi N,\phi P$ etc. Every number coprime to the product $A$ is also coprime to each of its factors $M,N,P$ etc., and conversely (Article 19). Also, every number congruent to one of $m,m',m''$ etc. modulo $M$ is coprime to $M$, and conversely; the same holds for $N,P$ etc. So the question becomes: how many numbers below $A$ are congruent, modulo $M$, to one of $m,m',m''$ etc., modulo $N$, to one of $n,n',n''$
\SourceLeaf{36}
etc., and so on? By Article 32, all numbers giving specified residues modulo each of $M,N,P$ etc. are congruent modulo their product $A$. So below $A$ there is just one number congruent to given residues modulo each of $M,N,P$ etc. The number we seek is therefore the number of combinations of each of $m,m',m''$ with each of $n,n',n''$ and each of $p,p',p''$ etc. The theory of combinations shows that this number is $=\phi M.\,\phi N.\,\phi P$ etc. \mbox{\emph{Q. E. D.}}

IV. It is now easy to see how to apply this to the case in question. Resolve $A$ into its prime factors, or write it in the form $a^\alpha b^\beta c^\gamma$ etc., where $a,b,c$ etc. are distinct primes. Then
\[
\begin{aligned}
\phi A&=\phi a^\alpha.\,\phi b^\beta.\,\phi c^\gamma\ \mathrm{etc.}\\
&=a^{\alpha-1}(a-1)\,b^{\beta-1}(b-1)\,c^{\gamma-1}(c-1)\ \mathrm{etc.}
\end{aligned}
\]
or, more neatly,
\[\phi A=A\,\frac{a-1}{a}.\,\frac{b-1}{b}.\,\frac{c-1}{c}\ \mathrm{etc.}\]

\emph{Example.} Let $A=60=2^2.\,3.\,5$, so that $\phi A=\frac{1}{2}.\,\frac{2}{3}.\,\frac{4}{5}.\,60=16$. The numbers coprime to $60$ are $1,7,11,13,17,19,23,29,31,37,41,43,47,49,53.\,59.$

The first solution of this problem appears in the distinguished Euler's paper \emph{theoremata arithmetica nova methodo demonstrata}, Comm. nov. Ac. Petrop. VIII, p. 74. The proof was later repeated in another paper, \emph{Speculationes circa quasdam insignes proprietates numerorum}, Acta Petrop. VIII, p. 17.

\Article{39}
If we define the symbol $\phi$ so that $\phi A$ gives the number of numbers coprime to $A$ and \emph{not greater than} $A$, then clearly $\phi1$ is no longer $=0$, but $=1$. All other cases remain unchanged. Adopting this definition, we have the following theorem.

\emph{If $a,a',a''$ etc. are all the divisors of $A$, including unity and $A$ itself, then}
\[\phi a+\phi a'+\phi a''+\textit{etc.}=A\]

\emph{Example.} Let $A=30$; then $\phi1+\phi2+\phi3+\phi5+\phi6+\phi10+\phi15+\phi30=1+1+2+4+2+4+8+8=30$.

\emph{Proof.} Multiply all the numbers coprime to $a$ and not greater than $a$ by $\frac{A}{a}$, and likewise all those coprime to $a'$ by $\frac{A}{a'}$, and so on. This gives $\phi a+\phi a'
+\phi a''+$ etc. numbers, none greater than $A$. Now
\SourceLeaf{37}

1) these numbers are all distinct. Those generated from the \emph{same} divisor of $A$ are clearly distinct. If equal numbers were generated from different divisors $M,N$ and numbers $\mu,\nu$ coprime to them, respectively, that is, if $\frac{A}{M}\mu=\frac{A}{N}\nu$, then $\mu N=\nu M$. We may assume $M>N$. Since $M$ is coprime to $\mu$ and divides $\mu N$, it also divides $N$: the greater divides the smaller. \emph{Q. E. A.}

2) These numbers include every one of $1,2,3....A$. Let any number not exceeding $A$ be $t$, and let the greatest common divisor of $A,t$ be $\delta$. Then $\frac{A}{\delta}$ is a divisor of $A$ coprime to $\frac{t}{\delta}$. Clearly the number $t$ is among those generated from the divisor $\frac{A}{\delta}$.

3) So the number of these numbers is $A$, and therefore
\[\phi a+\phi a'+\phi a''+\mathrm{etc.}=A.\quad\emph{Q. E. D.}\]

\Article{40}
\emph{If the greatest common divisor of $A,B,C,D$ etc. is $=\mu$, numbers $a,b,c,d$ etc. can be found such that}
\[aA+bB+cC+\textit{etc.}=\mu\]

\emph{Proof.} First consider just two numbers $A,B$, with greatest common divisor $=\lambda$. The congruence $Ax\equiv\lambda\ (\mathrm{mod.}\ B)$ is then solvable (Article 30). Let a root be $\equiv\alpha$, and put $\frac{\lambda-A\alpha}{B}=\beta$. Then $\alpha A+\beta B=\lambda$, as required.

On adding a third number $C$, let the greatest common divisor of the numbers $\lambda,C$ be $=\lambda'$. This is also the greatest common divisor of $A,B,C$\HistoricalNote{*)}{For $\lambda'$ clearly divides all of $A,B,C$. If it were not the \emph{greatest} common divisor, that divisor would be greater than $\lambda'$. But, since this greatest common divisor divides $A,B,C$, it also divides $k\alpha A+k\beta B+\gamma C$, that is, $\lambda'$: the greater divides the smaller. \emph{Q. E. A.} --- This follows even more easily from Article 18.}. Find numbers $k,\gamma$ such that $k\lambda+\gamma C=\lambda'$; then $k\alpha A+k\beta B+\gamma C=\lambda'$.

On adding a fourth number $D$, let the greatest common divisor of $\lambda',D$, which is easily seen to be also the greatest common divisor of $A,B,C,D$, be $=\lambda''$, and put $k'\lambda'+\delta D=\lambda''$. Then $kk'\alpha A+kk'\beta B+k'\gamma C+\delta D=\lambda''$.

\SourceLeaf{38}
We can proceed in the same way however many other numbers are added.

So, if the numbers $A,B,C,D$ etc. have no common divisor, we can clearly obtain
\[aA+bB+cC+\mathrm{etc.}=1\]

\Article{41}
\emph{If $p$ is prime and there are $p$ objects, any number of which may be equal, provided they are not all equal, then the number of permutations of these objects is divisible by $p$.}

\emph{Example.} The five objects $A,A,A,B,B$ can be arranged in ten different ways.

This theorem is easily proved from the familiar theory of permutations. Suppose there are $a$ equal objects $=A$, then $b$ equal objects $=B$, then $c$ equal objects $=C$, and so on, where any of the numbers $a,b,c$ etc. may also be unity, so that
\[a+b+c+\mathrm{etc.}=p\]
The number of permutations is
\[=\frac{1.2.3......p}{1.2.3...a.1.2..b.1.2..c\ \mathrm{etc.}}\]
The numerator of this fraction is clearly divisible by the denominator, since the number of permutations must be an integer. The numerator is divisible by $p$, whereas the denominator, being composed of factors smaller than $p$, is not divisible by $p$ (Article 15). So the number of permutations is divisible by $p$ (Article 19).

We hope, though, that some readers will also enjoy the following proof.

Call two permutations \emph{similar} when the order of their objects differs only in this way: the object occupying the first position in one occupies another position in the other, while the remaining objects follow in the same order, and the object that was first follows the one that was last\HistoricalNote{*)}{If we imagine similar permutations written around a circle, with the last object next to the first, they do not differ at all, since no position can be called first or last.}. Thus, in our example, the permutations $ABAAB$ and $ABABA$ are similar, since the objects in the first, second, etc. positions of the former occur in the third, fourth, etc. positions of the latter, in the same order.

\SourceLeaf{39}

Each permutation contains $p$ objects, so we can obtain $p-1$ similar permutations by moving the object that was first to the second, third, etc. positions. If no two of these can be identical, then the total number of permutations is clearly divisible by $p$, since it is $p$ times the number of mutually dissimilar permutations. So suppose that two permutations

\[PQ\ldots TV\ldots YZ;\qquad V\ldots YZPQ\ldots T\]

\noindent obtained from each other by shifting the terms are identical, that is, $P=V$ etc. Let the term $P$, which is first in the former, occupy position $n+1$ in the latter. Then, in the latter sequence, the term in position $n+1$ equals the first, that in position $n+2$ equals the second, and so on. Thus the term in position $2n+1$ again equals the first, and the same holds for $3n+1$ etc. In general, the term in position $kn+m$ equals that in position $m$. Here, if $kn+m$ exceeds $p$, either the sequence $V\ldots YZPQ\ldots T$ is to be repeated from the beginning, or we subtract from $kn+m$ the nearest smaller multiple of $p$. If we now choose $k$ so that $kn\equiv1\ (\mathrm{mod.}p)$, which is possible because $p$ is prime, it follows in general that the term in position $m$ equals that in position $m+1$. Thus every term equals the next, \emph{that is}, all terms are equal, contrary to the hypothesis.

\Article{42}

\emph{If the coefficients $A,B,C\ldots.N$, $a,b,c\ldots.n$ of two functions of the form}

\[
\begin{array}{l@{\qquad}c}
x^m+Ax^{m-1}+Bx^{m-2}+Cx^{m-3}\mathinner{\ldotp\ldotp\ldotp\ldotp\ldotp}+N
\quad\mathinner{\ldotp\quad\ldotp\quad\ldotp\quad\ldotp\quad\ldotp}&(P)\\
x^\mu+ax^{\mu-1}+bx^{\mu-2}+cx^{\mu-3}\mathinner{\ldotp\ldotp\ldotp\ldotp\ldotp}+n
\quad\mathinner{\ldotp\quad\ldotp\quad\ldotp\quad\ldotp\quad\ldotp}&(Q)
\end{array}
\]

\noindent\emph{are all rational, but not all integers, and the product of $(P)$ and $(Q)$ is}

\[=x^{m+\mu}+\mathfrak{A}x^{m+\mu-1}+\mathfrak{B}x^{m+\mu-2}+\emph{etc.}+\mathfrak{Z}\]

\noindent\emph{then its coefficients $\mathfrak{A},\mathfrak{B}\ldots.\mathfrak{Z}$ cannot all be integers.}

\emph{Proof.} Express all fractions among the coefficients $A,B$ etc., $a,b$ etc. in lowest terms, and choose any prime $p$ dividing one or more of their denominators. We may assume that $p$ divides the denominator of a fractional coefficient in $(P)$. If $(Q)$ is divided by $p$, then $\frac{(Q)}{p}$ clearly also has at least one fractional coefficient whose denominator contains the factor $p$, namely the first coefficient $\frac{1}{p}$.

\SourceLeaf{40}

There is clearly a fractional term in $(P)$ whose denominator contains a \emph{higher} power of $p$ than the denominators of all preceding fractional terms, and a power \emph{no lower} than the denominators of all following terms. Let this term be $=Gx^g$, and let the number of factors of $p$ in the denominator of $G$ be $=t$. There is a similar term in $\frac{(Q)}{p}$; let it be $=\varGamma x^\gamma$, and let the number of factors of $p$ in the denominator of $\varGamma$ be $=\tau$. Clearly $t+\tau$ is at least $=2$. With these preparations, the term $x^{g+\gamma}$ in the product of $(P)$ and $(Q)$ has a fractional coefficient whose denominator contains $t+\tau-1$ factors of $p$. The proof is as follows.

In $(P)$, let the terms preceding $Gx^g$ be ${}'Gx^{g+1}$, ${}''Gx^{g+2}$ etc., and those following it $G'x^{g-1}$, $G''x^{g-2}$ etc. Similarly, in $\frac{(Q)}{p}$, let $\varGamma x^\gamma$ be preceded by ${}'\varGamma x^{\gamma+1}$, ${}''\varGamma x^{\gamma+2}$ etc. and followed by $\varGamma'x^{\gamma-1}$, $\varGamma''x^{\gamma-2}$ etc. Then, in the product of $(P)$, $\frac{(Q)}{p}$, the coefficient of $x^{g+\gamma}$ is

\[
\begin{aligned}
&=G\varGamma+{}'G\varGamma'+{}''G\varGamma''+\text{etc.}\\
&\phantom{=G\varGamma}\mathord{+}{}'\varGamma G'+{}''\varGamma G''+\text{etc.}
\end{aligned}
\]

The term $G\varGamma$ is a fraction whose denominator, in lowest terms, contains $t+\tau$ factors of $p$. If fractional, the other terms have fewer factors of $p$ in their denominators, since each is a product of two factors, one involving no more than $t$ and the other fewer than $\tau$ factors of $p$, or one no more than $\tau$ and the other fewer than $t$. Thus $G\varGamma$ has the form $\frac{e}{fp^{t+\tau}}$, and the sum of the remaining terms has the form $\frac{e'}{f'p^{t+\tau-\delta}}$, where $\delta$ is positive and $e,f,f'$ are free of the factor $p$. Their total sum is therefore $=\frac{ef'+e'fp^\delta}{ff'p^{t+\tau}}$. Its numerator is not divisible by $p$, so no reduction can leave fewer than $t+\tau$ such factors in the denominator. The coefficient of $x^{g+\gamma}$ in the product of $(P),(Q)$ is consequently

\[=\frac{ef'+e'fp^\delta}{ff'p^{t+\tau-1}}.\]

\noindent\emph{that is, a fraction} whose denominator contains $t+\tau-1$ factors of $p$.\par
\noindent\mbox{\emph{Q. E. D.}}

\SourceLeaf{41}

\Article{43}

\emph{A congruence of degree $m$}

\[Ax^m+Bx^{m-1}+Cx^{m-2}+\emph{etc.}+Mx+N\equiv0\]

\noindent\emph{whose modulus is a prime $p$ not dividing $A$ cannot be solved in more than $m$ different ways, or have more than $m$ roots incongruent modulo $p$} (see Articles 25, 26).

If anyone denies this, suppose that there are congruences of various degrees $m,n$ etc. having more than $m,n$ etc. roots, and let the smallest such degree be $m$, so that all such congruences of lower degrees satisfy our theorem. Since it has already been proved for the first degree (Article 26), $m$ must be $=2$ or greater. The congruence

\[Ax^m+Bx^{m-1}+\text{etc.}+Mx+N\equiv0\]

\noindent therefore has at least $m+1$ roots, say $x\equiv\alpha,x\equiv\beta,x\equiv\gamma$ etc. We may suppose all the numbers $\alpha,\beta,\gamma$ etc. positive and smaller than $p$, with $\alpha$ the least. Substitute for $x$ the expression $y+\alpha$, so that the given congruence becomes

\[A'y^m+B'y^{m-1}+C'y^{m-2}....+M'y+N'\equiv0\]

Clearly this is satisfied by setting $y\equiv0$, or $\equiv\beta-\alpha$, or $\equiv\gamma-\alpha$, etc. All these roots are distinct, and their number is $=m+1$. Since $y\equiv0$ is a root, $N'$ must be divisible by $p$. So the expression

\[y(A'y^{m-1}+B'y^{m-2}+\text{etc.}+M')\text{ is }\equiv0\ (\mathrm{mod.}p),\]

\noindent when $y$ is assigned one of the $m$ values $\beta-\alpha,\gamma-\alpha$ etc., all of which are $>0$ and $<p$. In each of these cases, therefore, we also have

\[A'y^{m-1}+B'y^{m-2}+\text{etc.}+M'\text{ is }\equiv0\ \text{(Article 22)}\]

\noindent\emph{that is}, the congruence \qquad $A'y^{m-1}+B'y^{m-2}+\text{etc.}+M'\equiv0$

\noindent of degree $m-1$ has $m$ roots and so contradicts our theorem. For $A'$ is clearly $=A$ and therefore not divisible by $p$, as required. Yet we assumed that every congruence of degree lower than $m$ satisfies the theorem. \mbox{\emph{Q. E. A.}}

\SourceLeaf{42}

\Article{44}

Although we assumed here that the modulus $p$ does not divide the coefficient of the highest term, the theorem is not restricted to this case. If the first coefficient, or even some of the following ones, were divisible by $p$, these terms could be discarded, eventually lowering the degree of the congruence until its first coefficient was no longer divisible by $p$. This holds unless all the coefficients are divisible by $p$, in which case the congruence is an identity and the unknown is entirely unrestricted.

The distinguished La Grange was the first to state and prove this theorem (\emph{M\'em. de l'Ac. de Berlin, Ann\'ee} 1768, \emph{p.} 192). It also appears in the distinguished Le Gendre's paper \emph{Recherches d'Analyse ind\'etermin\'ee, Hist. de l'Acad. de Paris} 1785, \emph{p.} 466. The distinguished Euler proved in \emph{Nov. Comm. Ac. Petr.} XVIII, \emph{p.} 93 that the congruence $x^n-1\equiv0$ cannot have more than $n$ distinct roots. Although this is a special case, the method used by this great man is easily adapted to all congruences. He had previously dealt with an even more restricted case in \emph{Comm. nov. Ac. Petr.} V, \emph{p.} 6, but that method cannot be applied generally. In Section VIII below we will prove the theorem in yet another way. However different all these methods may appear at first, those familiar with the subject who compare them will easily see that they all rest on the same principle. Since we consider this theorem here only as a lemma, and a full treatment does not belong here, we will not discuss composite moduli separately.

\BookRule{0.28\linewidth}
